\documentclass[11pt]{article}
\usepackage[a4paper,margin=20mm]{geometry}
\usepackage{amsmath,amssymb,xcolor,graphicx}
\usepackage[hypcap=false]{caption}
\usepackage[hidelinks,bookmarksnumbered,bookmarksopen]{hyperref}
\hypersetup{pdftitle={Quantum mechanics 4 — The particle in a box},pdfauthor={Lecturer: Prof. R. Okafor},
  pdfcreator={KherveNote}}
\renewcommand\labelitemii{\textbullet}
\renewcommand\labelitemiii{\textbullet}
\renewcommand\labelitemiv{\textbullet}
\renewcommand\labelenumii{\arabic{enumi}.\arabic{enumii}.}
\renewcommand\labelenumiii{\arabic{enumi}.\arabic{enumii}.\arabic{enumiii}.}
\renewcommand\labelenumiv{\arabic{enumi}.\arabic{enumii}.\arabic{enumiii}.\arabic{enumiv}.}
\setlength{\parindent}{0pt}
\setlength{\parskip}{0.6em}
\definecolor{knoteaccent}{HTML}{1A6DD8}
\definecolor{knotemuted}{HTML}{6B6B6B}
\definecolor{knotekey}{HTML}{D96B00}
\newcommand\knotetime[1]{\leavevmode\llap{\color{knotemuted}\scriptsize #1\hspace{1em}}}
\newcommand\knotekey[2][]{\par\noindent#1\colorbox{knotekey!12}{\parbox{\dimexpr\linewidth-2\fboxsep}{%
  \textbf{\color{knotekey}$\star$ Key point.} #2}}\par}
\newcommand\knotequestion[2][]{\par\noindent#1{\color{knoteaccent}\textbf{?}}~\emph{#2}\par}
\newenvironment{knotetranscript}{\par\color{knotemuted}\small}{\par}
\newcommand\knotesummary[1]{\par\noindent\fcolorbox{knoteaccent}{knoteaccent!6}{%
  \parbox{\dimexpr\linewidth-2\fboxsep-2\fboxrule}{\textbf{Summary}\par #1}}\par}
\newenvironment{knotechunk}{}{}
\pagestyle{plain}
\begin{document}
\begin{knotechunk}
{\color{knoteaccent}\rule{\linewidth}{1.5pt}}\par
{\LARGE\bfseries Quantum mechanics 4 — The particle in a box}\par
{\large Lecturer: Prof. R. Okafor}\hfill{\color{knotemuted}06 October 2026 \textperiodcentered{} Physics lecture theatre 1}\par
{\color{knoteaccent}\rule{\linewidth}{0.6pt}}\par
\knotesummary{A particle confined to a box of width L can only have energies E\textsubscript{n} = n\textsuperscript{2}h\textsuperscript{2}/(8mL\textsuperscript{2}): energy is quantised, the lowest level is not zero, and the levels spread apart as the box shrinks — the origin of quantum confinement in nanoparticles and quantum dots.}

\section{Setting up the problem}

\knotetime{01:55}Potential $V = 0$ for $0 < x < L$ and $V = \infty$ outside: the particle cannot leave the box. Inside, the time-independent Schrödinger equation is \[-\frac{\hbar^2}{2m}\frac{d^2\psi}{dx^2} = E\psi\] with boundary conditions $\psi(0) = \psi(L) = 0$.
\end{knotechunk}

\begin{knotechunk}
\section{Solutions}

\begin{itemize}\setlength{\itemsep}{0pt}\setlength{\parskip}{0pt}
  \item general solution $\psi = A\sin kx + B\cos kx$, with $k = \sqrt{2mE}/\hbar$
  \item $\psi(0) = 0 \Rightarrow B = 0$
  \item $\psi(L) = 0 \Rightarrow kL = n\pi$, $n = 1, 2, 3, \dots$
\end{itemize}

\knotetime{09:40}Normalised wavefunctions and energies: \[\psi_n(x) = \sqrt{\tfrac{2}{L}}\,\sin\frac{n\pi x}{L}, \qquad E_n = \frac{n^2\pi^2\hbar^2}{2mL^2} = \frac{n^2 h^2}{8 m L^2}\]

\knotekey[\knotetime{10:05}]{Energy is quantised by the boundary conditions alone, and $E_1 > 0$: zero-point energy, consistent with the uncertainty principle.}

\begin{itemize}\setlength{\itemsep}{0pt}\setlength{\parskip}{0pt}
  \item $\psi_n$ has $n - 1$ nodes inside the box
  \item level spacing $E_{n+1} - E_n \propto 2n + 1$ grows with $n$
  \item probability density $|\psi_n|^2$ is not uniform; for large $n$ it averages to $1/L$ (correspondence principle)
\end{itemize}
\end{knotechunk}

\begin{knotechunk}
\section{Numbers: an electron in a 1 nm box}

\knotetime{20:25}$E_1 = h^2/(8 m_e L^2) = (6.626\times10^{-34})^2 / (8 \times 9.109\times10^{-31} \times 10^{-18}) \approx 6.0\times10^{-20}$ J $\approx 0.38$ eV.

\begin{itemize}\setlength{\itemsep}{0pt}\setlength{\parskip}{0pt}
  \item $E_2 = 4E_1 \approx 1.5$ eV, so the $1 \to 2$ transition is about 1.1 eV (in the near infrared)
  \item for a 1 g marble in a 1 cm box the spacing is immeasurably small — classical behaviour
\end{itemize}
\end{knotechunk}

\begin{knotechunk}
\section{Quantum confinement in materials}

\knotetime{29:25}$E \propto 1/L^2$: the smaller the box, the larger the gap.

\begin{itemize}\setlength{\itemsep}{0pt}\setlength{\parskip}{0pt}
  \item quantum dots (e.g. CdSe): smaller dots emit bluer light
  \item quantum wells in semiconductor lasers; conjugated molecules (free-electron model of polyenes)
\end{itemize}

\knotequestion[\knotetime{30:40}]{Why does a finite well always have at least one bound state, while the infinite one starts at n = 1?}
\end{knotechunk}

\begin{knotechunk}
\section{What was said}
\begin{knotetranscript}
\knotetime{11:00:50}Today we solve our first real problem, a particle trapped in a box with infinitely high walls.

\knotetime{11:06:20}Inside the box the potential is zero, so the Schrödinger equation is just the free particle one.

\knotetime{11:09:40}The wavefunction has to vanish at both walls, and that's what forces the energy to come in steps.

\knotetime{11:14:30}Notice the lowest energy is not zero. You cannot have a particle at rest in a box, that would violate uncertainty.

\knotetime{11:18:20}Let's put numbers in for an electron in a box one nanometre wide.

\knotetime{11:21:40}That gives about zero point four electronvolts for the ground state, so these are energies you can see with light.

\knotetime{11:27:10}Now make the box smaller and the levels spread apart as one over L squared.

\knotetime{11:30:30}This is exactly why small cadmium selenide quantum dots glow blue and larger ones red.
\end{knotetranscript}
\end{knotechunk}
\end{document}
